\chapter{The Existing SR5E1 CAN Library}
\label{ch:arch}

This chapter describes the library as it was found: the offset convention
that maps Motor Control codes onto CAN identifiers, the periodic diagnostic
stack, the transmission and reception paths, and the command map. Defects
are noted where they affect the reading of the code; the repairs themselves
are postponed to \cref{ch:dev}.

\section{Architectural overview}
\label{sec:arch-overview}

The library is an adapter. On one side is the dynamometer software, which
speaks a flat set of 11-bit CAN identifiers. On the other side is the ST
Motor Control SDK, which speaks the Frame Communication Protocol recalled
in \cref{sec:mcp}. Between them sits \texttt{DYNO2DWProtocol}, a translation
function invoked from the FDCAN receive interrupt.

\begin{figure}[ht]
\centering
\begin{tikzpicture}[
  font=\sffamily\small,
  box/.style={draw=NavyInk, rounded corners=1.5pt, align=center,
              inner sep=7pt, minimum height=12mm, minimum width=32mm},
  arr/.style={-{Latex}, thick, NavyInk}
]
  \node[box, fill=Paper] (dyno) {Dyno bench\\software};
  \node[box, fill=Paper, right=18mm of dyno] (xlate)
    {DYNO2DW\\{\scriptsize identifier $\leftrightarrow$ FCP}};
  \node[box, fill=Paper, right=18mm of xlate] (fw)
    {Stellar-E firmware\\{\scriptsize ST MC protocol}};
  \draw[arr] (dyno) -- node[above, font=\scriptsize] {CAN $+$\texttt{0x30}}
                        node[below, font=\scriptsize] {commands} (xlate);
  \draw[arr] (xlate) -- node[above, font=\scriptsize] {Code, Len,}\\
                        node[below, font=\scriptsize] {Payload, CRC} (fw);
  \draw[arr] ([yshift=-7mm]fw.south) -- ++(0,-7mm) -|
             node[near start, below, font=\scriptsize] {diagnostics $+$\texttt{0x40}}
             ([yshift=-7mm]dyno.south);
\end{tikzpicture}
\caption{Role of the DYNO2DW translation layer. Commands travel dyno
$\rightarrow$ MCU with an identifier offset of $+$\texttt{0x30}; diagnostics
travel MCU $\rightarrow$ dyno with an offset of $+$\texttt{0x40}.}
\label{fig:arch}
\end{figure}

The adapter is deliberately thin. It does not implement a second motor
controller and it does not keep a shadow copy of the PI gains. It unpacks
a CAN frame, builds an FCP frame, and hands it to
\texttt{CFCP\_RX\_IRQ\_Handler}. Status in the other direction is read with
\texttt{UI\_GetReg} and emitted as raw 4-byte little-endian values under
the diagnostic identifiers.

\section{The offset rule}
\label{sec:offset}

Two constants structure the identifier map.

\begin{definition}[Command offset]
An incoming command whose Motor Control code is $c$ is published on the
bus as identifier $c + \mathtt{0x30}$.
\end{definition}

\begin{definition}[Diagnostic offset]
An outgoing diagnostic whose Motor Control register identifier is $r$ is
published on the bus as identifier $r + \mathtt{0x40}$.
\end{definition}

The constants are not arbitrary. They park the command window in
\texttt{0x30}--\texttt{0x3F} and the diagnostic window in
\texttt{0x40}--\texttt{0x60}, leaving \texttt{0x10}--\texttt{0x11} free for
MCU responses (ACK/NACK and \texttt{REG\_VALUE}) and avoiding collision
with the very lowest identifiers, which many toolchains reserve. The same
rule is what later allowed unused command slots such as \texttt{0x3B} and
\texttt{0x3C} to be claimed for independent $I_d$/$I_q$ writes without
renumbering the rest of the map.

\begin{remark}
The offset is applied to the \emph{protocol code}, not to an arbitrary
enum. When a new Motor Control code is introduced---for example
\texttt{MC\_PROTOCOL\_CODE\_SET\_CURRENT\_ID\_REF = 0x0F}---its CAN
identifier is implied: $\mathtt{0x0F}+\mathtt{0x30}=\mathtt{0x3F}$. In
practice a collision check against the diagnostic window and against
already assigned commands is still required; that is why the DBC chapter
records \texttt{0x3B}/\texttt{0x3C} as the identifiers actually shipped
for the independent current commands, reserving the ``naive'' offset
values where they would have overlapped existing ramps.
\end{remark}

\section{Periodic diagnostics: \texttt{Send\_Status}}
\label{sec:send-status}

The dyno must observe the inverter without polling. The library therefore
emits a stack of ten diagnostic frames every \SI{100}{\milli\second}, from
a function named \texttt{Send\_Status}. The registers in the stack are:

\begin{center}
\begin{tabular}{cll}
\toprule
Index & Register macro & Physical meaning \\
\midrule
0 & \texttt{MC\_PROTOCOL\_REG\_UNDEFINED} & Sequence counter \\
1 & \texttt{MC\_PROTOCOL\_REG\_SPEED\_REF} & Speed reference \\
2 & \texttt{MC\_PROTOCOL\_REG\_TORQUE\_REF} & $I_q$ reference \\
3 & \texttt{MC\_PROTOCOL\_REG\_FLUX\_REF} & $I_d$ reference \\
4 & \texttt{MC\_PROTOCOL\_REG\_BUS\_VOLTAGE} & DC-link voltage \\
5 & \texttt{MC\_PROTOCOL\_REG\_HEATS\_TEMP} & Heatsink temperature \\
6 & \texttt{MC\_PROTOCOL\_REG\_SPEED\_MEAS} & Measured speed \\
7 & \texttt{MC\_PROTOCOL\_REG\_TORQUE\_MEAS} & Measured $I_q$ \\
8 & \texttt{MC\_PROTOCOL\_REG\_FLUX\_MEAS} & Measured $I_d$ \\
9 & \texttt{MC\_PROTOCOL\_REG\_STATUS} & Machine state \\
\bottomrule
\end{tabular}
\end{center}

Each iteration of the loop reads one register through
\texttt{UI\_GetReg}, packs the 32-bit value as four little-endian bytes,
assigns the identifier $r + \mathtt{0x40}$, and submits the frame with
\texttt{DBCCAN\_TX\_IRQ\_Handler}. Submission is gated on
\texttt{FCP\_TRANSFER\_IDLE}: if the TX path is still busy from the
previous frame, the entire stack is skipped for that period. That
back-pressure rule is conservative---it prefers silence over a torn
update---and it is one of the reasons the validation in \cref{ch:val}
starts by confirming that ten identifiers really do appear, in order,
every \SI{100}{\milli\second}.

The resulting diagnostic identifiers, after the offset, are collected in
\cref{tab:diag-map}. They are the same identifiers the GUI plots in
\cref{ch:gui}.

\begin{table}[ht]
\centering
\caption{Outgoing diagnostic identifiers (offset $+$\texttt{0x40}).}
\label{tab:diag-map}
\begin{tabular}{lll}
\toprule
ID & Name & Content \\
\midrule
\texttt{0x42} & \texttt{DIAG\_STATE\_MC}     & Machine status \\
\texttt{0x44} & \texttt{DIAG\_SPEED\_REF}    & Speed reference \\
\texttt{0x48} & \texttt{DIAG\_TORQUE\_REF}   & $I_q$ reference \\
\texttt{0x4C} & \texttt{DIAG\_FLUX\_REF}     & $I_d$ reference \\
\texttt{0x59} & \texttt{DIAG\_BUS\_VOLTAGE}  & DC-link voltage \\
\texttt{0x5A} & \texttt{DIAG\_TEMPERATURE}   & Heatsink temperature \\
\texttt{0x5E} & \texttt{DIAG\_SPEED\_MEAS}   & Measured speed \\
\texttt{0x5F} & \texttt{DIAG\_TORQUE\_MEAS}  & Measured $I_q$ \\
\texttt{0x60} & \texttt{DIAG\_FLUX\_MEAS}    & Measured $I_d$ \\
\bottomrule
\end{tabular}
\end{table}

\section{Transmission path}
\label{sec:tx}

For a single diagnostic frame the TX path is four steps.

\begin{enumerate}[leftmargin=1.6em]
  \item \texttt{UI\_GetReg(\&pHandle->\_Super, Register, NULL)} returns a
        signed 32-bit value. For PI registers the value is already packed
        as described in \cref{sec:pi-pack}; for physical quantities it is
        the SDK's internal representation (RPM, \si{\milli\ampere},
        \si{\milli\volt}, \si{\celsius}, \ldots).
  \item The four data bytes are written little-endian into the FDCAN TX
        buffer. Little-endian matches both the Cortex-R52 memory order and
        the convention used by Vector tools on Intel hosts, so a DBC signal
        of type \texttt{unsigned Motorola} would be wrong here.
  \item The identifier is set to the register identifier plus
        \texttt{0x40}.
  \item \texttt{DBCCAN\_TX\_IRQ\_Handler} submits the buffer to the FDCAN
        peripheral. Hardware then performs arbitration and acknowledgement
        autonomously.
\end{enumerate}

Command responses---ACK/NACK and \texttt{GET\_REG} values---take a
slightly different path, through \texttt{CFCP\_TX\_IRQ\_Handler}. That
handler is where the original library used a single identifier for both
kinds of response; the defect and its repair are in \cref{sec:ack}.

\section{Reception path}
\label{sec:rx}

Reception is entirely interrupt-driven.

\begin{enumerate}[leftmargin=1.6em]
  \item Hardware raises \texttt{IRQ\_CAN1\_LINE0\_HANDLER} on a matching
        identifier.
  \item The ISR copies the payload into \texttt{RxArr[]} and calls
        \texttt{DYNO2DWProtocol}.
  \item The function switches on the command code implied by the
        identifier (\texttt{EXECUTE\_COMMAND}, \texttt{SET\_RAMP},
        \texttt{SET\_REG}, \ldots), writes \texttt{msg[0]} (FCP code),
        \texttt{msg[1]} (length) and the payload bytes, then appends
        \texttt{msg[n] = SelfCRCCalc(msg)}.
  \item Each byte of \texttt{msg[]} is fed to
        \texttt{CFCP\_RX\_IRQ\_Handler}, which advances the FCP state
        machine and, on a complete well-checksummed frame, dispatches into
        the Motor Control UI layer.
\end{enumerate}

Timeouts live in the FCP layer, not in DYNO2DW: an interrupted multi-byte
feed is abandoned so that a later command cannot be parsed as the tail of
an earlier one. This is the correct behaviour on a bus that may drop a
frame, but it makes the ISR path sensitive to the order and completeness
of the \texttt{msg[]} construction. Several of the bugs repaired in
\cref{ch:dev} were, at bottom, bugs in that construction---a length byte
that did not match the payload, a command that wrote \texttt{msg[]} twice,
a CRC computed before the payload was filled.

\section{Command map}
\label{sec:cmd-map}

\Cref{tab:cmd-map} lists the incoming commands as documented at the
beginning of the project. Identifiers follow the $+$\texttt{0x30} rule
except where a later collision forced a reassignment; those exceptions are
discussed with the corresponding feature in \cref{ch:dev}.

\begin{table}[ht]
\centering
\caption{Incoming command identifiers (offset $+$\texttt{0x30}).}
\label{tab:cmd-map}
\begin{tabular}{llp{70mm}}
\toprule
ID & Command & Role \\
\midrule
\texttt{0x31} & \texttt{SET\_REG}            & Write a Motor Control register \\
\texttt{0x32} & \texttt{GET\_REG}            & Read a Motor Control register \\
\texttt{0x33} & \texttt{EXECUTE\_COMMAND}    & Start, stop, reset, align, \ldots \\
\texttt{0x37} & \texttt{SET\_RAMP}           & Speed ramp (final value, duration) \\
\texttt{0x3A} & \texttt{SET\_CURRENT\_ID\_IQ}& Simultaneous $I_d$ and $I_q$ \\
\texttt{0x3B} & \texttt{SET\_CURRENT\_ID}    & Independent $I_d$ reference \\
\texttt{0x3C} & \texttt{SET\_CURRENT\_IQ}    & Independent $I_q$ reference \\
\texttt{0x3E} & \texttt{SET\_ID\_RAMP}       & Flux ($I_d$) ramp \\
\texttt{0x3F} & \texttt{SET\_IQ\_RAMP}       & Torque ($I_q$) ramp \\
\bottomrule
\end{tabular}
\end{table}

\texttt{EXECUTE\_COMMAND} is a bit-mask, not a single enumerated action.
Byte~0 carries \texttt{EXE\_START\_MOTOR}, \texttt{STOP\_RAMP},
\texttt{RESET}, \texttt{EXE\_PING}, \texttt{FAULT\_ACK} and
\texttt{ENCODER\_ALIGN}; byte~1 carries \texttt{EXE\_IQDREF\_CLEAR} and
\texttt{GET\_BOARD\_INFO}. The original decoder treated those flags with a
chain of independent \texttt{if} statements, which is the first defect
repaired in the next chapter.

\section{What the original library did not provide}
\label{sec:gaps}

A short list, so that Chapter~\ref{ch:dev} can be read as a closing of
named gaps rather than as a tour of unrelated patches.

\begin{itemize}[leftmargin=1.6em]
  \item \textbf{Unambiguous machine commands.} Start and stop could be
        requested together; the last write to \texttt{msg[2]} won, and a
        second FCP frame could still be emitted for reset.
  \item \textbf{Independent current axes.} $I_d$ and $I_q$ could be written
        as a pair; they could not be written, or ramped, separately with
        first-class protocol codes.
  \item \textbf{A flux ramp.} The torque ramp existed in the SDK
        (\texttt{UI\_ExecTorqueRamp}). The flux ramp did not; it had to be
        added through the UI, MCI and STC layers.
  \item \textbf{Atomic PI updates.} Numerator and divisor travelled in
        ways that a single \texttt{SET\_REG} could not express.
  \item \textbf{Distinguishable responses.} ACK/NACK and
        \texttt{GET\_REG} shared an identifier.
\end{itemize}

Those five items are the body of \cref{ch:dev}.
